\section{Coarse information outside a sublinear collar}

The first information estimate of \cite[Section~3]{OpenAI2026PolynomialPEPS}
bounds the mutual information of the lattice sample $T$ of a bounded union of
rectangles by reading $T$ as a template on the cut $\Lambda$ and applying the
sublinear collar bound (\cite[Proposition~9.6]{OpenAI2026AreaLaw}) and its
amplification.

\begin{definition}[Bounded unions of lattice boxes]
    \label{def:peps_box_union}
    \lean{TNLean.PEPS.Approximation.IsBoxUnion}
    \leanok
    A finite set $T\subseteq\mathbb Z^2$ is a union of at most $N_0$ boxes if
    $T=\bigcup_{a<k}[\ell_a,h_a]$ for some $k\le N_0$ and $\ell_a,h_a\in\mathbb Z^2$,
    where $[\ell,h]:=\{x\in\mathbb Z^2:\ell_i\le x_i\le h_i\ (i=1,2)\}$ (possibly empty).
\end{definition}

\begin{lemma}[Lattice samples of rectangles]
    \label{lem:peps_box_union_of_intervals}
    \lean{TNLean.PEPS.Approximation.isBoxUnion_of_ordConnected}
    \leanok
    \uses{def:peps_box_union}
    For real intervals $I_a,J_a\subseteq\mathbb R$ ($a<k\le N_0$, open, closed, or half-open), any finite set $T=\mathbb Z^2\cap\bigcup_{a<k}(I_a\times J_a)$ is a union of at most $N_0$ boxes (Definition~\ref{def:peps_box_union}).
\end{lemma}
\begin{proof}
    \leanok
    Since $T$ is finite, each nonempty factor $(I_a\cap\mathbb Z)\times(J_a\cap\mathbb Z)$ equals $[\ell_a,h_a]$ with $\ell_a=(\min(I_a\cap\mathbb Z),\min(J_a\cap\mathbb Z))$ and $h_a=(\max(I_a\cap\mathbb Z),\max(J_a\cap\mathbb Z))$, while an empty factor is replaced by an empty box.
\end{proof}

\begin{lemma}[Translations and templates]
    \label{lem:peps_box_union_template}
    \lean{TNLean.PEPS.Approximation.IsBoxUnion.translate,
        TNLean.PEPS.Approximation.IsBoxUnion.exists_template}
    \leanok
    \uses{def:peps_box_union,def:area_law_template}
    If $T$ is a union of at most $N_0$ boxes, so is $T+d$ for any $d\in\mathbb Z^2$. If $T\ne\varnothing$ has $\operatorname{diam}_\infty T\le D$ with $D+1\le s_0$, $D\le n$, and $\max(C_{\rm tpl},0)N_0(s_0+1)\le n$, then $T$ is an $(n,s_0)$-template (Definition~\ref{def:area_law_template}) with closed-rectangle pieces.
\end{lemma}
\begin{proof}
    \leanok
    Each nonempty box $[\ell,h]$ is $\mathbb Z^2\cap[\ell-(\tfrac14,\tfrac14),h+(\tfrac14,\tfrac14)]$ (of $\ell_\infty$-diameter at most $D+\tfrac12<s_0$), and an empty box is $\mathbb Z^2\cap[-\tfrac14,-\tfrac18]^2$.
\end{proof}

\begin{lemma}[Coarse information outside a sublinear collar]
    \label{thm:peps_coarse_information}
    \lean{TNLean.PEPS.Approximation.coarse_mutualInformation_le_square,
        TNLean.PEPS.Approximation.coarse_mutualInformation_le_translated_square}
    \leanok
    \uses{def:peps_box_union,def:polynomial_peps_model,
        thm:peps_coarse_information_of_collar,thm:area_law_sublinear_collar}
    In \cite[Proposition~3.1]{OpenAI2026PolynomialPEPS}, for every integer
    $N_0\ge1$ and $A_0\ge1$ there exist $\beta<1$ and $C,u_0<\infty$ such that for every $u\ge u_0$, every lattice sample $T=\mathbb Z^2\cap\bigcup_{a<k}R_a$ of $k\le N_0$ axis-parallel rectangles with $\operatorname{diam}_\infty T\le A_0u$, every $X\subseteq\Lambda\cap T$, and every $E\subseteq\Lambda$ with $\operatorname{dist}_\infty(E,T)>u^\beta$,
    \begin{align}
        I_\Omega(X:E) & \le C,
        \label{eq:peps_coarse_information}
    \end{align}
    including for $X=\Lambda\cap T$, consistently open or closed rectangle sides, and after any common coordinate translation.
\end{lemma}
\begin{proof}
    \leanok
    \uses{def:area_law_collar_bound,thm:area_law_sublinear_collar,
        thm:peps_coarse_information_of_collar,
        thm:peps_coarse_information_square,lem:peps_box_union_of_intervals,
        lem:peps_box_union_template}
    Lemma~\ref{lem:peps_box_union_of_intervals} shows that $T$ (or its pullback by a translation $d$) is a union of at most $N_0$ boxes, so combining Lemma~\ref{thm:peps_coarse_information_square} with Lemma~\ref{thm:area_law_sublinear_collar} (\cite[Proposition~9.6]{OpenAI2026AreaLaw}) at $D_0=13$ and $C_{\rm tpl}=24$ gives (\ref{eq:peps_coarse_information}).
\end{proof}

\begin{lemma}[The coarse exponent is below one]
    \label{lem:peps_coarse_exponent_lt_one}
    \lean{TNLean.PEPS.Approximation.beta_lt_one}
    \leanok
    The amplification exponent $\beta_0$ satisfies $\beta_0<1$.
\end{lemma}

\begin{proof}
    \leanok
    Direct numerical evaluation of $\beta_0$.
\end{proof}

% Source: OpenAI2026PolynomialPEPS, proof of Proposition 3.1 (prop:coarse),
% 02-information.tex, lines 110--146.
\begin{definition}[Constants of the coarse information estimate]
    \label{def:peps_coarse_information_constants}
    \lean{TNLean.PEPS.Approximation.coarseScaleFactor,
        TNLean.PEPS.Approximation.coarseInformationThreshold}
    \leanok
    \uses{thm:area_law_uniform_collar_threshold}
    For $C_{\rm tpl},A_0\in\mathbb R$, $N_0\in\mathbb Z$, collar threshold $N$, and amplification constants $(C_a,N_a)$, let $N_r$ be the threshold of Lemma~\ref{thm:area_law_uniform_collar_threshold} for $C_a$ and $\eta=K_{\rm c}^{-\beta_0}$ (so $\lfloor n^{1-\varepsilon}\rfloor+2r<K_{\rm c}^{-\beta_0}n^{\beta_0}$ for $n\ge N_r$ and $0\le r\le C_an^{(1-\varepsilon)/\alpha}$), where
    \begin{align}
        K_{\rm c}(C_{\rm tpl},N_0,A_0)     & =(\max\{C_{\rm tpl},0\}+1)(N_0+1)
        (\max\{A_0,1\}+4), \notag                                              \\
        u_0(C_{\rm tpl},N_0,A_0,N,C_a,N_a) & =
        \max\{N,N_a,N_r,2,K_{\rm c}^{1/\varepsilon}\}.
        \label{eq:peps_coarse_information_constants}
    \end{align}
\end{definition}

\begin{lemma}[Coarse information with explicit constants]
    \label{thm:peps_coarse_information_of_collar_explicit}
    \lean{TNLean.PEPS.Approximation.coarse_mutualInformation_le_of_sublinearCollarBoundAt}
    \leanok
    \uses{def:peps_box_union,def:area_law_collar_bound,
        def:area_law_amplification_bound,def:peps_coarse_information_constants}
    Suppose the sublinear collar bound (Definition~\ref{def:area_law_collar_bound}) holds at $(q,R,J,\Delta,D_0,C_{\rm tpl})$ with $(C,N)$, and the amplification bound (Definition~\ref{def:area_law_amplification_bound}) holds at $(\max\{C,1\},1)$ with $(C_a,N_a)$. For $u\ge u_0(C_{\rm tpl},N_0,A_0,N,C_a,N_a)$ from (\ref{eq:peps_coarse_information_constants}), any unit ground vector $\Omega$ with gap $\Delta$ of an admissible range-$R$, bound-$J$ Hamiltonian on a finite induced domain $\Lambda$, any union $T$ of at most $N_0$ boxes with $\operatorname{diam}_\infty T\le A_0u$, any $X\subseteq\Lambda\cap T$, and any $E\subseteq\Lambda$ with $\operatorname{dist}_\infty(E,T)>u^{\beta_0}$,
    \begin{align}
        I_\Omega(X:E) & \le \max\{C_a,0\}. \notag
    \end{align}
\end{lemma}
\begin{proof}
    \leanok
    \uses{lem:peps_box_union_template,def:peps_coarse_information_constants,
        thm:area_law_uniform_collar_threshold}
    For $T\ne\varnothing$, set $A=\max(A_0,1)$, $c_1=(\max(C_{\rm tpl},0)+1)(N_0+1)$, $s_0=\lceil Au\rceil+1$, and $n=\lceil c_1(Au+3)\rceil$. Since $u\ge u_0\ge\max\{2,K_{\rm c}^{1/\varepsilon}\}$, Lemma~\ref{lem:peps_box_union_template} makes $T$ an $(n,s_0)$-template separated from the empty boundary of the cut $\Lambda$, with $u\le n\le K_{\rm c}u$ and $n^{1-\varepsilon}\le K_{\rm c}^{1-\varepsilon}u^{1-\varepsilon}\le u\le s_0$. Applying the collar bound and the amplification bound at $(\max(C,1),1)$ with $n\ge u\ge N_r$ gives $\lfloor n^{1-\varepsilon}\rfloor+2r<K_{\rm c}^{-\beta_0}n^{\beta_0}\le u^{\beta_0}$ by Lemma~\ref{thm:area_law_uniform_collar_threshold}, so $E\cap N_{2r}(S_0)=\varnothing$ and $I_\Omega(X:E)\le C_an^{-1}\le\max(C_a,0)$.
\end{proof}

\begin{lemma}[Coarse information, assuming the sublinear collar bound]
    \label{thm:peps_coarse_information_of_collar}
    \lean{TNLean.PEPS.Approximation.coarse_mutualInformation_le_of_sublinearCollarBound}
    \leanok
    \uses{def:peps_box_union,def:area_law_collar_bound,
        thm:area_law_amplification,lem:peps_box_union_template,
        thm:area_law_uniform_collar_threshold}
    Fix $q,R,J$, $\Delta>0$, $N_0$, and $A_0$, and assume the sublinear collar bound (Definition~\ref{def:area_law_collar_bound}) for some $D_0,C_{\rm tpl}$. Then there exist $\beta<1$, $C$, and $u_0$ such that for every $u\ge u_0$, unit ground vector $\Omega$ with gap $\Delta$ of an admissible range-$R$, bound-$J$ Hamiltonian on a finite induced domain $\Lambda$, union $T$ of at most $N_0$ boxes with $\operatorname{diam}_\infty T\le A_0u$, $X\subseteq\Lambda\cap T$, and $E\subseteq\Lambda$ with $\operatorname{dist}_\infty(E,T)>u^\beta$, one has $I_\Omega(X:E)\le C$.
\end{lemma}
\begin{proof}
    \leanok
    \uses{thm:peps_coarse_information_of_collar_explicit,thm:area_law_amplification,
        def:peps_coarse_information_constants,lem:peps_coarse_exponent_lt_one}
    Combine the assumed collar constants $(C,N)$ with the amplification constants $(C_a,N_a)$ of Theorem~\ref{thm:area_law_amplification} at $(\max\{C,1\},1)$, Lemma~\ref{lem:peps_coarse_exponent_lt_one} ($\beta=\beta_0<1$), and Lemma~\ref{thm:peps_coarse_information_of_collar_explicit}.
\end{proof}

\begin{lemma}[Coarse information on the square grid, assuming the collar bound]
    \label{thm:peps_coarse_information_square}
    \lean{TNLean.PEPS.Approximation.coarse_mutualInformation_le_square_of_sublinearCollarBound,
        TNLean.PEPS.Approximation.finiteProductMutualInformation_relabel}
    \leanok
    \uses{thm:peps_coarse_information_of_collar,def:polynomial_peps_model,
        lem:peps_box_union_template}
    Fix $q$, $J\in\mathbb R$, $\Delta>0$, $N_0$, and $A_0$, and assume the range-$1$ sublinear collar bound. There exist $\beta<1$, $C$, and $u_0$ such that for every $u\ge u_0$, open square grid with a model Hamiltonian (Definition~\ref{def:polynomial_peps_model}), unit ground vector $\Omega$ with gap $\Delta$, translation $c\in\mathbb Z^2$, union $T$ of at most $N_0$ boxes with $\operatorname{diam}_\infty T\le A_0u$, and site sets $X,E$ with $X+c\subseteq T$ and $\operatorname{dist}_\infty(E+c,T)>u^\beta$, one has $I_\Omega(X:E)\le C$.
\end{lemma}
\begin{proof}
    \leanok
    Translating $T$ by $-c$ preserves box-union structure (Lemma~\ref{lem:peps_box_union_template}), diameters, distances, and regional entropies, reducing the claim to Lemma~\ref{thm:peps_coarse_information_of_collar} at $R=1$.
\end{proof}

\begin{lemma}[Coarse information on finite induced domains]
    \label{thm:peps_coarse_information_finite_domain}
    \lean{TNLean.PEPS.Approximation.coarse_mutualInformation_le}
    \leanok
    \uses{def:peps_box_union,def:area_law_hamiltonian}
    For integers $q\ge1$, $R,N_0\ge0$ and reals $J\ge0$, $\Delta>0$, $A_0$, there exist $\beta<1$, $C$, and $u_0$ such that for every $u\ge u_0$, finite induced domain $\Lambda$, unit ground vector $\Omega$ with gap $\Delta$ of an admissible range-$R$, bound-$J$ Hamiltonian on $\Lambda$, union $T$ of at most $N_0$ boxes with $\operatorname{diam}_\infty T\le A_0u$, $X\subseteq\Lambda\cap T$, and $E\subseteq\Lambda$ with $\operatorname{dist}_\infty(T,E)>u^\beta$, one has $I_\Omega(X:E)\le C$.
\end{lemma}
\begin{proof}
    \leanok
    \uses{thm:peps_coarse_information_of_collar,thm:area_law_sublinear_collar}
    Apply Lemma~\ref{thm:area_law_sublinear_collar} at $D_0=2R+11$ and $C_{\rm tpl}=24$, followed by Lemma~\ref{thm:peps_coarse_information_of_collar}.
\end{proof}

\begin{definition}[Closed-form collar constants at range $R$]
    \label{def:peps_closed_collar_constants}
    \lean{TNLean.PEPS.Approximation.collarBoxConstantClosed,
        TNLean.PEPS.Approximation.collarConstantClosed,
        TNLean.PEPS.Approximation.collarThresholdClosed}
    \leanok
    \uses{def:area_law_collar_bound}
    For $q,R\in\mathbb N$ and $J,\Delta\in\mathbb R$, let $C_{\rm box}(q,R,J,\Delta)$ be the safe-box constant of \cite[Proposition~9.5]{OpenAI2026AreaLaw} at range $R$ with transport constants $(\text{threshold }1/8,\ \text{entropy }200000\text{ of degree }2,\ \text{energy }3(18\pi^2+6144e^{\pi^2+1})\text{ of degree }4)$, and let $C(q,R,J,\Delta),N(q,R,J,\Delta)$ be the resulting collar constant and threshold of \cite[Proposition~9.6]{OpenAI2026AreaLaw} at truncation factor $10$.
\end{definition}

\begin{lemma}[The sublinear collar bound at range $R$ at closed-form constants]
    \label{thm:peps_collar_closed}
    \lean{TNLean.PEPS.Approximation.sublinearCollarBoundAt_closed}
    \leanok
    \uses{def:area_law_collar_bound,def:peps_closed_collar_constants}
    For $q,R\in\mathbb N$ with $q\ge1$, $J\ge0$, $\Delta>0$, and $D_0>2R+10$, the sublinear collar bound (Definition~\ref{def:area_law_collar_bound}) at $(R,D_0,C_{\rm tpl}=24)$ holds with $C(q,R,J,\Delta)$ and $N(q,R,J,\Delta)$ of Definition~\ref{def:peps_closed_collar_constants}.
\end{lemma}
\begin{proof}
    \leanok
    \uses{thm:area_law_sublinear_collar_explicit,
        thm:area_law_transport_bound_explicit}
    Combine Lemma~\ref{thm:area_law_transport_bound_explicit} with Lemma~\ref{thm:area_law_sublinear_collar_explicit} at $C_{\rm tpl}=24$.
\end{proof}

\begin{definition}[Closed-form coarse constants at range $R$]
    \label{def:peps_closed_coarse_constants}
    \lean{TNLean.PEPS.Approximation.coarseInformationConstantClosed,
        TNLean.PEPS.Approximation.coarseInformationThresholdClosed}
    \leanok
    \uses{def:peps_closed_collar_constants,
        def:peps_collar_amplification_constants,
        def:peps_coarse_information_constants}
    For $q,R,N_0\in\mathbb N$ and $J,\Delta,A_0\in\mathbb R$, with $C,N$ from Definition~\ref{def:peps_closed_collar_constants} and $C_a=C_a(q,R,J,\Delta,C)$, $N_a=N_a(q,R,J,\Delta,C)$ from Definition~\ref{def:peps_collar_amplification_constants}, define via (\ref{eq:peps_coarse_information_constants}):
    \begin{align}
        C_{\rm coarse}(q,R,J,\Delta) & := \max\{C_a,0\},                    &
        u_0(q,R,J,\Delta,N_0,A_0)    & := u_0(24,N_0,A_0,N,C_a,N_a). \notag
    \end{align}
\end{definition}

% Status: auxiliary result. Proposition 3.1 of the polynomial-PEPS source
% asserts only that constants exist, on the original square at range one; that
% statement is thm:peps_coarse_information_square.
% Source: polynomial-PEPS 02-information.tex, Proposition 3.1, lines 91--145.
% Revision: adc7f1241b42e322a6451854ab7e4b4c146bf78a.
\begin{lemma}[Coarse information at range $R$ at closed-form constants]
    \label{thm:peps_coarse_information_closed}
    \lean{TNLean.PEPS.Approximation.coarse_mutualInformation_le_closed}
    \leanok
    \uses{def:peps_box_union,def:area_law_hamiltonian,
        def:peps_closed_coarse_constants}
    For $q,R,N_0\in\mathbb N$ ($q\ge1$), $J\ge0$, $\Delta>0$, $A_0\in\mathbb R$, $u\ge u_0(q,R,J,\Delta,N_0,A_0)$, any unit ground vector $\Omega$ with gap $\Delta$ of an admissible range-$R$, bound-$J$ Hamiltonian on a finite induced domain $\Lambda$, any union $T$ of at most $N_0$ boxes with $\operatorname{diam}_\infty T\le A_0u$, any $X\subseteq\Lambda\cap T$, and any $E\subseteq\Lambda$ with $\operatorname{dist}_\infty(T,E)>u^{\beta_0}$, one has $I_\Omega(X:E)\le C_{\rm coarse}(q,R,J,\Delta)$ (extending \cite[Proposition~3.1]{OpenAI2026PolynomialPEPS} from $R=1$ to general $R$).
\end{lemma}
\begin{proof}
    \leanok
    \uses{thm:peps_coarse_information_of_collar_explicit,thm:peps_collar_closed,
        thm:peps_collar_amplification_bound}
    Combine Lemma~\ref{thm:peps_collar_closed} (at $D_0=2R+11$) and Lemma~\ref{thm:peps_collar_amplification_bound} with Lemma~\ref{thm:peps_coarse_information_of_collar_explicit}.
\end{proof}
